\documentclass[11pt]{amsart}
\usepackage[T1]{fontenc}
\usepackage{lmodern}
\usepackage{amsmath,amssymb,amsthm,booktabs}
\usepackage[margin=1in]{geometry}
\usepackage[colorlinks=true,linkcolor=blue,citecolor=blue,urlcolor=blue]{hyperref}
\hypersetup{pdftitle={Balanced random subset sums in finite abelian groups: explicit corrections of order sqrt(q log q)},pdfauthor={Haoyu Chen}}
\allowdisplaybreaks
\newtheorem{theorem}{Theorem}
\newtheorem{lemma}{Lemma}
\newtheorem*{lemma2prime}{Lemma 2$'$}
\newtheorem*{propM}{Proposition M}
\newtheorem*{corC}{Corollary C}
\newtheorem*{proposition}{Proposition}
\theoremstyle{remark}
\newtheorem*{remark}{Remark}
\newcommand{\Z}{\mathbb Z}
\newcommand{\Q}{\mathbb Q}
\newcommand{\F}{\mathbb F}
\newcommand{\E}{\mathbb E}
\newcommand{\PI}{\mathbb P_{\mathrm I}}
\newcommand{\PU}{\mathbb P_{\mathrm U}}
\newcommand{\Fail}{\mathcal F}
\newcommand{\Hom}{\operatorname{Hom}}
\newcommand{\ord}{\operatorname{ord}}
\newcommand{\lo}{P_{\mathrm{low}}}
\newcommand{\hi}{P_{\mathrm{high}}}
\newcommand{\LB}{\Lambda_B}
\newcommand{\LC}{\Lambda_C}
\newcommand{\QB}{Q_B}
\title[Explicit corrections for balanced subset sums]{Balanced random subset sums in finite abelian groups:\ explicit corrections of order $\sqrt{q\log q}$}
\author{Haoyu Chen}
\address{Independent researcher}
\date{12 October 2026}
\subjclass[2020]{Primary 11B13; Secondary 20K01, 60C05}
\keywords{Random subset sums, finite abelian groups, character sums, integer lattices}
\begin{document}
\begin{abstract}
Let $G$ be any finite abelian group of order $N$, and put $q=\log_2N$.
We give explicit sufficient sample sizes for all subset-sum representation
counts to lie within relative error $1/2$ of $2^k/N$. For independent
uniform samples, every integer $k\ge q+8400\sqrt{q\log_2q}$ has failure
probability at most $1/(2q)$ when $N\ge2^{64}$. The constant is $200$
when $N\ge2^{486018750}$. For a uniform $k$-element subset the failure
bound is $1/q$, provided also $k\le\sqrt N$. We give two intermediate
explicit bounds and a corollary with $k=\lceil q+3q^{4/5}\rceil$ at the
larger onset. The proof splits a Fourier power according to integer
relation lattices, bounds the low-rank part deterministically on a
common event, and takes an unconditional second moment of the remaining
polynomial. A count of bounded-generator lattices yields the square-root
correction. The leading-coefficient-one question is already known to
have an affirmative answer; the subject here is its explicit correction,
uniformly over all finite abelian groups.
\end{abstract}
\maketitle

\section{Introduction and statements}\label{sec:intro}
Write the operation of a finite abelian group $G$ additively, and let
$N=|G|$. For a tuple $A=(g_1,\ldots,g_k)$ define
\[
 R_A(g)=\#\bigl\{S\subseteq\{1,\ldots,k\}:\ \sum_{i\in S}g_i=g\bigr\}.
\]
The empty subset has sum zero. Repeated entries retain distinct indices.
In model \textup{(I)}, the entries are independent and uniform in $G$.
In model \textup{(U)}, $A$ is a uniform $k$-element subset of $G$, and
$R_A$ counts subsets of this set. In either model put
\[
 q=\log_2N,\qquad L=\log_2q,\qquad
 D(g)=\frac{NR_A(g)}{2^k}-1,\qquad
 \Fail=\{\max_{g\in G}|D(g)|>1/2\}.
\]
Thus ``balanced'' means the non-strict inequality
$\max_g|D(g)|\le1/2$. The numbers $q$ and $L$ are real, and are never
assumed integral. Every probability assertion below is for each fixed
group separately, with the same bound for every group of the given order.

The question recorded as Erd\H{o}s Problem~1179 \cite{EP1179} asks for
leading coefficient one in the sample size for near-equal subset-sum
counts, with high probability. Erd\H{o}s and Hall \cite{EH76} already
settled that question. We study explicit smaller corrections. We retain
the following two intermediate results to display the successive
parameter choices; Theorem~\ref{thm:sqrt} is the main estimate.

\begin{theorem}\label{thm:log}
Let $G$ be any finite abelian group of order $N\ge2^{2^{20}}$, and let
\[
 k=\left\lceil q+\frac{1000q}{\log_2q}\right\rceil.
\]
Then $\PI(\Fail)\le1/(2q)$. Moreover $k\le N$, and
$\PU(\Fail)\le1/q$.
\end{theorem}

\begin{theorem}\label{thm:power}
Let $G$ be any finite abelian group of order $N\ge2^{64}$, and let
$k=\lceil q+3400q^{4/5}\rceil$. Then $\PI(\Fail)\le1/(2q)$.
Moreover $k\le N$, and $\PU(\Fail)\le1/q$.
\end{theorem}

\begin{theorem}\label{thm:sqrt}
Let $G$ be any finite abelian group of order $N$.
\begin{enumerate}
\renewcommand{\labelenumi}{(\alph{enumi})}
\item If $N\ge2^{64}$, every integer $k\ge q+8400\sqrt{qL}$ satisfies
$\PI(\Fail)\le1/(2q)$. If also $k\le2^{q/2}$, then
$\PU(\Fail)\le1/q$.
\item If $N\ge2^{486018750}$, the same assertions hold with $200$ in
place of $8400$, retaining the condition $k\le2^{q/2}$ for model
\textup{(U)}.
\end{enumerate}
The admissible integer interval in model \textup{(U)} is nonempty in
both parts.
\end{theorem}

\begin{corC}\label{cor:power}
For every finite abelian group of order $N\ge2^{486018750}$ and
$k=\lceil q+3q^{4/5}\rceil$, one has $\PI(\Fail)\le1/(2q)$ and
$\PU(\Fail)\le1/q$.
\end{corC}

There are no asymptotic error terms in these statements or their proofs.
For the minimal rounded size in Theorem~\ref{thm:sqrt}, the correction
is at most $C\sqrt{qL}+1$, with $C=8400$ or $200$ and exactly the onsets
displayed above. In particular, using $L\le4\sqrt q$ for $q\ge1$,
\begin{equation}\label{eq:explicit-asymptotic}
 \frac{\lceil q+C\sqrt{qL}\rceil-q}{q}
 \le2Cq^{-1/4}+q^{-1},\qquad
 \frac{\sqrt{qL}}{q^{4/5}}\le2q^{-1/20}.
\end{equation}
The second ratio is at most $\epsilon>0$ for
$q\ge\max\{1,(2/\epsilon)^{20}\}$; the first is at most $\epsilon$
for $q\ge\max\{1,(4C/\epsilon)^4,2/\epsilon\}$.
These quantify the asymptotic language in the title.

\begin{remark}[Finite-size scope]\label{rem:scope}
The elementary $L^2$ estimate already gives the tuple failure bound for
$k\ge2q+L+3$. The theorems above do not improve this sample requirement
throughout all their stated ranges. Here are exact sufficient comparison
ranges, with all comparisons made at the \emph{minimal} allowed integer
$k$. In Theorem~\ref{thm:log},
\[
 \frac{1000q}{L}\le q\quad\Longleftrightarrow\quad
 L\ge1000\quad\Longleftrightarrow\quad q\ge2^{1000}.
\]
Thus for $q\ge2^{1000}$ (equivalently $N\ge2^{2^{1000}}$) its size is
at most $\lceil2q\rceil<2q+L+3$. In Theorem~\ref{thm:power} the
corresponding condition is
\[
 q\ge3400^5=454\,354\,240\,000\,000\,000,
\]
and for Theorem~\ref{thm:sqrt}(a) it is
$q/L\ge8400^2=70\,560\,000$. Below each of these respective
conditions the unrounded size, and hence its ceiling, exceeds $2q$;
any improvement over $\lceil2q+L+3\rceil$ is then less than $L+4$,
if there is an improvement at all. These are not claimed to be the
exact crossover points after the logarithmic term and ceilings are
included. For a correction $c(q)$ the exact integer comparison is
\[
 \lceil q+c(q)\rceil<\lceil2q+L+3\rceil
 \quad\Longleftrightarrow\quad
 q+c(q)\le\lceil2q+L+3\rceil-1.
\]
Theorem~\ref{thm:sqrt}(b) improves the elementary size throughout its
range: with $\QB=486\,018\,750$, monotonicity of $q/\log_2q$ gives
$q/L\ge\QB/\log_2\QB>\QB/29>40000$, so $200\sqrt{qL}<q$.
At the onset of Theorem~\ref{thm:log}, its size is $53\,477\,376$,
whereas the elementary size is $2\,097\,175$. Constants have not been
optimized. The transfer to model \textup{(U)} is proved separately in
Section~\ref{sec:transfer}.
\end{remark}

\section{Related work}\label{sec:related}
Erd\H{o}s and R\'enyi \cite[Theorem 1, pp.~131--132]{ER65} give
simultaneous relative error at most $\epsilon$ with probability greater
than $1-\delta$ for independent uniform samples if
\[
 k\ge2\log_2N+2\log_2(1/\epsilon)+\log_2(1/\delta).
\]
With $\epsilon=1/2$ and $\delta=1/(2q)$ this is precisely the
$2q+L+3$ benchmark. Erd\H{o}s and Hall
\cite[unnumbered Theorem, p.~174]{EH76} prove simultaneous
$(1\pm\eta)2^k/N$ counts for almost all independent uniform tuples
in every finite abelian group, for fixed $\eta>0$, when
\[
 k>q\{1+O_\eta(\log\log q/\log q)\}.
\]
This cited historical statement has an unspecified implied constant
and no numerical finite onset or failure function. It is not used as
an effective premise here. As we read their proof, their Lemma~2, p.~176, drops a
product-principal constraint on characters and overcounts binary
relations; their Lemma~3 passes to a maximum-count estimate. Our proof
instead separates tuples before taking an unconditional second moment.

Hall's 1977 paper \cite{Hall77} is bibliographically confirmed, but its
theorem statements were unavailable in the literature search underlying
this comparison. We assign no mathematical conclusion to it. This
access gap, and incomplete forward-citation coverage, limit any claim
about the record.

Stronger results for restricted groups must also be distinguished.
For $G=(\Z/2)^d$, the representation function is constant on the
span of the sampled elements and zero off it. Erd\H{o}s and Hall
\cite[p.~449, Theorem 3 and preceding paragraph]{EH78} give failure
of coverage less than $2^{d-k}$; coverage is exact uniformity in this
family. Thus $k\ge d+\lceil\log_2(2d)\rceil$ suffices for the tuple
failure target here.

Dhar and Dvir \cite[Theorem 2.1]{DD24} establish simultaneous two-sided
relative guarantees for finite-field linear hashing. In the prime-field
case, for $S\subseteq\F_p^n$, $n\ge5$, $\tau>0$, $0<\delta<1$,
\[
 p\ge32\max\{n(1+\tau)/(\tau\delta)^2,n\},\qquad
 p^r<|S|\le p^{r+1},\qquad4\le r\le n-1,
\]
at least a $1-\delta$ fraction of surjective linear maps to
$\F_p^{r-3}$ give relative error at most $\tau$ at every point.
Their Theorems~2.2 and~2.4 give binary variants with logarithmic
entropy loss under explicit size hypotheses. These are genuine
pointwise balance results, with a finite-field domain.

Lund's preprint \cite[Corollary 1.4]{Lund26} gives a further two-sided
hashing statement for nonempty $S\subseteq\F_p^n$, $1\le d\le n$,
$0<\tau\le1$, and $0<\delta<1$. If an integer $a\ge1$ satisfies
\[
 p^{a-1}\ge16(d+a),\qquad
 \frac{|S|}{p^d}\ge\frac{192p^a}{\tau^2\delta},
\]
the relative error of every fiber is at most $\tau$ with probability
greater than $1-\delta$, for either uniform surjections or unrestricted
uniform matrices. Taking $S=\{0,1\}^k$, $N=p^d$, $\tau=1/2$, and
$\delta=1/(2q)$ makes the second condition
$2^{k-q}\ge1536qp^a$. For each fixed prime this gives logarithmic
overhead. We cite the preprint's theorem statement without claiming
an independent review of its proof. Neither the field-dependent
condition nor the domain supplies a bound uniform over all finite
abelian groups, and no without-replacement conclusion is inferred
without a transfer argument.

Our quantifier is over \emph{every} finite abelian group, with explicit
constants, onsets, and probabilities. Within a literature search of
12 October 2026 we did not locate an explicit bound of this form for
all finite abelian groups. This is a bounded search statement, with
Hall 1977 unread, and makes no priority claim. The paper does not claim
a new solution of the original leading-coefficient question.

\section{Proof of Theorem 1}\label{sec:proof1}
We supply the algebraic ingredients and the Fourier argument in full.
Write $H=\widehat G$ additively, with identity $0$, while evaluating
characters multiplicatively: $(\chi+\psi)(x)=\chi(x)\psi(x)$.
The structure theorem gives $|H|=N$. For $\Lambda\le\Z^s$, define
\[
 d_p(\Lambda)=\dim_{\F_p}\operatorname{im}(\Lambda\longrightarrow\F_p^s),
 \quad r_p(\Lambda)=s-d_p(\Lambda),\quad
 r_G(\Lambda)=\max_{p\mid N}r_p(\Lambda).
\]
Here $r_p(\Lambda)=\dim_{\F_p}((\Z^s/\Lambda)/p(\Z^s/\Lambda))$.
We use ``lattice'' for any subgroup of $\Z^s$, not necessarily full rank
or saturated.

\begin{lemma}\label{lem:hom}
For every $\Lambda\le\Z^s$,
\begin{equation}\label{eq:hom}
 |\Hom(\Z^s/\Lambda,H)|\le N^{r_G(\Lambda)}.
\end{equation}
If finitely generated abelian groups satisfy $S\le K$ and $S\twoheadrightarrow Q$,
then for every prime $p$,
\begin{equation}\label{eq:subquot}
 \dim(Q/pQ)\le\dim(S/pS)\le\dim(K/pK).
\end{equation}
\end{lemma}
\begin{proof}
Write $K$ as a free abelian group of rank $m$ plus finite torsion. If
its $p$-primary part has $b_p$ cyclic factors, then
$\dim(K/pK)=m+b_p$. Each free generator and each $p$-primary cyclic
generator has at most $|H_p|$ images in the $p$-primary part $H_p$ of
$H$. Torsion at other primes maps to zero, so
$|\Hom(K,H_p)|\le|H_p|^{m+b_p}$. Taking $K=\Z^s/\Lambda$ and
multiplying over $p\mid N$ proves \eqref{eq:hom}. A subgroup $S$
has free rank at most $m$, and $S[p]$ embeds in $K[p]$, so its number
of $p$-primary cyclic factors is at most $b_p$. This proves the second
inequality of \eqref{eq:subquot}. The first follows from the induced
surjection $S/pS\twoheadrightarrow Q/pQ$.
\end{proof}

\begin{lemma}\label{lem:coarse}
Let $s,B\ge1$ be integers. A lattice generated by a prescribed set of
vectors in $[-B,B]^s\cap\Z^s$ is generated by at most
\[
 J(s,B)=s+\lfloor\log_2(s!B^s)\rfloor
\]
vectors from that set. Hence there are at most
$(2B+1)^{sJ(s,B)}$ such lattices. For lattices generated by binary
vectors the number is at most
\begin{equation}\label{eq:binarypatterns}
 2^{s^2(1+\lceil\log_2s\rceil)}.
\end{equation}
\end{lemma}
\begin{proof}
Let $d$ be the rational rank and choose $d$ independent prescribed
generators, spanning $\Lambda_0$. Projection to $d$ coordinates with
nonsingular minor $M$ is injective on the rational span. It follows that
\[
 [\Lambda:\Lambda_0]\le|\det M|\le d!B^d\le s!B^s.
\]
The determinant estimate is the Leibniz expansion, with $d!$ terms
of size at most $B^d$. If the chosen vectors do not yet generate
$\Lambda$, adjoining one prescribed generator reduces the remaining
integer index by a factor at least two. At most
$\lfloor\log_2(s!B^s)\rfloor$ additions suffice. Rank zero is immediate.
Pad lists with zero vectors for counting; this gives the first count.
In the binary case
$J(s,1)\le s(1+\lceil\log_2s\rceil)$, and each entry has $2^s$
choices, which proves \eqref{eq:binarypatterns}.
\end{proof}

\begin{lemma}\label{lem:cube}
For every lattice $\Lambda\le\Z^s$, writing $r=r_G(\Lambda)$, one has
$|\Lambda\cap\{0,1\}^s|\le2^{s-r}$.
\end{lemma}
\begin{proof}
Choose $p\mid N$ attaining the maximum $r$. The image of $\Lambda$
modulo $p$ has dimension $s-r$. Some $s-r$ coordinate positions give
an injective projection on that image, by a nonsingular coordinate
minor of a basis. Integer binary vectors stay distinct modulo any prime,
including two. Their projections have at most $2^{s-r}$ binary choices.
\end{proof}

\begin{lemma}\label{lem:intersection}
If $t\ge1$ and $\Lambda\le\Z^{2t}$ is generated by binary vectors,
then $\Lambda\cap(\Z^t\times\{0\})$ is generated by vectors whose
first $t$ coordinates have absolute value at most
\begin{equation}\label{eq:box}
 B_t=2^{4t}t!.
\end{equation}
\end{lemma}
\begin{proof}
List distinct generators as columns of $(X;Y)$, with $t$ rows in each
block and $n\le2^{2t}$ columns. The intersection consists of $(Xz,0)$
for $z\in\Z^n$ with $Yz=0$. If $n=0$ it is zero. Let
$r=\operatorname{rank}_{\Q}Y$. If $r=0$, coordinate vectors generate
the kernel and their images have coordinates at most one. If $n=r$,
the kernel is zero. Otherwise choose $r$ columns and $r$ rows giving
a nonsingular matrix $C$, and write $F$ for the remaining columns on
those rows. The selected rows are independent, since $C$ is invertible,
and span the row space of $Y$, since its rank is $r$. With the pivot
coordinates first, the full kernel condition is therefore
\[
 z=(x,y),\qquad x=-C^{-1}Fy.
\]
Put $D_0=|\det C|$, so $1\le D_0\le r!$. The allowed integer $y$
form a lattice $T$ containing $D_0\Z^{n-r}$. Reducing coordinates
modulo $D_0$ shows that $T$ is generated by $D_0e_j$ and its elements
in $\{0,\ldots,D_0-1\}^{n-r}$. Every adjugate entry of $C$ has
absolute value at most $(r-1)!$. For a representative in this box,
each coordinate of $x$ is at most
\[
 \frac{r(r-1)!(n-r)(D_0-1)}{D_0}\le nr!.
\]
For $y=D_0e_j$ the bound is $r!$. The $y$ coordinates are also at
most $D_0\le r!$. Thus the integer kernel is generated by $z$ with
$\|z\|_\infty\le nr!$. Every coordinate of $Xz$ is at most
$n^2r!\le2^{4t}t!$, as required.
\end{proof}

For $u=(u_1,\ldots,u_t)\in H^t$ and $\theta\in H^{2t}$ set
\begin{align*}
 \LB(u)&=\left\langle c\in[-B_t,B_t]^t\cap\Z^t:
                         \sum_jc_ju_j=0\right\rangle,\\
 \LC(\theta)&=\left\langle e\in\{0,1\}^{2t}:
                         \sum_je_j\theta_j=0\right\rangle.
\end{align*}
\begin{lemma}\label{lem:persist}
If the first $t$ entries of $\theta$ are $u$, then
$r_G(\LB(u))\le r_G(\LC(\theta))$.
\end{lemma}
\begin{proof}
Identify the intersection of $\LC(\theta)$ with the first coordinate
block with a lattice $I\le\Z^t$. Every vector in $\LC(\theta)$ is
an actual relation. Lemma~\ref{lem:intersection} therefore generates
$I$ by bounded relations of $u$, and $I\le\LB(u)$. The coordinate
inclusion gives an injection
$\Z^t/I\hookrightarrow\Z^{2t}/\LC(\theta)$, and there is a
surjection $\Z^t/I\twoheadrightarrow\Z^t/\LB(u)$. Apply
\eqref{eq:subquot} for each $p\mid N$ and take maxima. The maximizing
primes need not coincide.
\end{proof}

\subsection{Fourier coefficients and a common event}
In model \textup{(I)}, let
\[
 a_\chi=\prod_{i=1}^k\frac{1+\chi(g_i)}2,\qquad h=k-q>0.
\]
Orthogonality gives
\begin{equation}\label{eq:fourier}
 D(g)=\sum_{\chi\ne0}a_\chi\overline{\chi(g)}.
\end{equation}
Indeed expanding the product and then summing all characters detects
each subset sum. A nonprincipal character has sum zero over $G$,
because translation by an element on which it is not one both preserves
and nontrivially multiplies that sum. All sums here are finite.

Let $\rho=2^{-q/16}$ and define
\[
 E=\{a_\chi=0\text{ for every }\ord(\chi)=2\}
 \cap\{|a_\chi|\le\rho\text{ for every }\ord(\chi)\ge3\}.
\]
An order-two coefficient is zero or one and is one with probability
$2^{-k}$. Its union bound is at most $N2^{-k}=2^{-h}$. If
$\ord(\chi)\ge3$, then for uniform $x\in G$
\[
 \E\left|\frac{1+\chi(x)}2\right|^4
 =\frac1{16}\E(6+4\chi(x)+4\chi(x)^{-1}+\chi(x)^2+\chi(x)^{-2})
 =\frac38.
\]
All four nonconstant characters in the expansion are nonprincipal.
Independence and Markov's inequality now give
\[
 \PI(\exists\chi,\ \ord(\chi)\ge3:\ |a_\chi|>\rho)
 \le N\rho^{-4}(3/8)^k
 \le2^{5q/4}(3/8)^q\le2^{-q/8}.
\]
The last step uses $\log_2(8/3)>11/8$, equivalently
$3^8=6561<8192=2^{13}$. Consequently
\begin{equation}\label{eq:E}
 \PI(E^c)\le2^{-h}+2^{-q/8}.
\end{equation}

\subsection{A split of the Fourier power}
Fix integers $t\ge6$, $1\le a\le t$, and let $\ell=\lceil\log_2t\rceil$.
Expand \eqref{eq:fourier} to obtain $D(g)^t=\lo(g)+\hi(g)$, where
$\lo$ collects the ordered nonprincipal tuples $u$ with
$r_G(\LB(u))<a$, and $\hi$ collects the rest. Multiplicities are
those of ordered tuples. Lemma~\ref{lem:coarse} bounds the number of
possible $\LB(u)$ by $2^{A_t}$, where
\begin{equation}\label{eq:oldA}
 A_t=2t^4(6+\ell)^2\le8t^6.
\end{equation}
To see every factor, use
\begin{align*}
 \log_2B_t&\le t(4+\ell),\\
 J(t,B_t)&\le t+t\ell+t^2(4+\ell)\le t^2(5+2\ell),\\
 t\log_2(2B_t+1)&\le t(\log_2B_t+2)\le t^2(6+\ell).
\end{align*}
The product is at most $2t^4(6+\ell)^2$. Since $\ell\le t$ and
$6+\ell\le2t$ for $t\ge6$, \eqref{eq:oldA} follows. For a fixed
low pattern $\Lambda$, the tuples are among the homomorphisms from
$\Z^t/\Lambda$ to $H$. Writing $r=r_G(\Lambda)<a$, their number is
at most $N^r\le N^{a-1}\le N^a$. Thus on $E$, for every $g$,
\begin{equation}\label{eq:low}
 |\lo(g)|\le2^{A_t+aq-tq/16}.
\end{equation}

The crucial second moment is unconditional:
\begin{equation}\label{eq:high}
 \E|\hi(g)|^2\le2^{C_t-ah},\qquad
 C_t=4t^2(2+\ell)\le8t^3.
\end{equation}
Indeed for two high tuples $u,v$, set $\theta=(u,-v)$. Independence
and orthogonality give the exact nonnegative expectation
\begin{equation}\label{eq:moment}
 \E\left[\prod_ja_{u_j}\prod_j\overline{a_{v_j}}\right]
 =\left[2^{-2t}\#\{e\in\{0,1\}^{2t}:\sum_je_j\theta_j=0\}\right]^k.
\end{equation}
The phase from evaluating at $g$ has absolute value one. The triangle
inequality therefore bounds the expanded second moment by the sum of
the right sides of \eqref{eq:moment}. Group pairs by
$\Lambda=\LC(\theta)$ and put $r=r_G(\Lambda)$. The binary vectors
of $\Lambda$ are exactly the binary relations of $\theta$: all
generators are relations and all binary relations were included.
Lemma~\ref{lem:cube} bounds the bracket by $2^{-r}$, and
Lemma~\ref{lem:persist} gives $r\ge a$. By \eqref{eq:hom} at most
$N^r$ tuples $\theta$, hence at most $N^r$ pairs $(u,v)$, have this
pattern. Its contribution is at most
$N^r2^{-rk}=2^{-rh}\le2^{-ah}$. Lemma~\ref{lem:coarse} gives at
most $2^{C_t}$ binary patterns, since
$\lceil\log_2(2t)\rceil=1+\lceil\log_2t\rceil$. Finally
$2+\ell\le2t$ for $t\ge2$, proving \eqref{eq:high}.
The event $E$ is absent from this expectation. No signed polynomial
is treated as a nonnegative random variable.

\subsection{The parameters for Theorem 1}
Assume first $L\ge500$ and take $t=\lfloor L\rfloor$,
$a=\lfloor t/32\rfloor$. For every integer $t\ge64$,
\begin{equation}\label{eq:induction}
 512t^5\le2^t.
\end{equation}
At $t=64$ the left side is $2^{39}\le2^{64}$. Its successive ratio
is at most $(65/64)^5<(9/8)^5<2$, the last comparison being
$59049<65536$. Induction proves \eqref{eq:induction}. Since
$q=2^L\ge2^t$, it follows that $A_t\le8t^6\le tq/64$, and
\[
 |\lo(g)|\le2^{-tq/64}\le2^{-t-1}\quad(g\in G)\quad\text{on }E.
\]
Here $a\le t/32$, $q\ge128$, and $2t\ge t+1$ justify the exponents.
Also
\[
 C_t+2t+2\le8t^3+2t+2\le16t^3\le512t^5\le q.
\]
As $t\ge500$, we have $a\ge t/32-1\ge t/40$, and
\[
 ah\ge\frac{t}{40}\frac{1000q}{L}
 \ge\frac{25tq}{t+1}\ge24q,
\]
where $L<t+1$ and $t\ge24$. If $E$ holds and $|D(g)|>1/2$,
then $|\hi(g)|>2^{-t-1}$. Markov's inequality for $|\hi(g)|^2$
and a union bound over $N$ points therefore give
\begin{equation}\label{eq:oldfailure}
 \PI(\Fail)\le2^{-h}+2^{-q/8}
       +2^{q+2t+2+C_t-ah}
 \le2^{-h}+2^{-q/8}+2^{-22q}.
\end{equation}
For real $x\ge4$, $2^x\ge x^2$: equality holds at four and the
logarithmic derivative is $\ln2-2/x>0$ thereafter, since
$\ln2>1/2$. Thus $q\ge L^2$. For $L\ge20$,
$L^2\ge8(L+3)$, by its value at twenty and positive derivative of
the difference. We obtain $h\ge1000q/L\ge1000L\ge L+3$ and
$q/8\ge L+3$. Each term in \eqref{eq:oldfailure} is consequently
at most $1/(8q)$, so the sum is at most $3/(8q)$.

For $20\le L\le500$ the stated size gives $k\ge3q$. Orthogonality
gives $\E|a_\chi|^2=2^{-k}$ for every nonprincipal character. Parseval
in \eqref{eq:fourier}, followed by Markov's inequality, gives the
general elementary estimate
\begin{equation}\label{eq:L2}
 \E\sum_g|D(g)|^2=N(N-1)2^{-k},\qquad
 \PI(\Fail)\le4\,2^{2q-k}.
\end{equation}
In this range it is at most $4\,2^{-q}\le1/(2q)$, since
$q\ge2^{20}$ and $2^q\ge q^2\ge8q$. The two ranges cover all
$L\ge20$, including their common boundary. This proves the tuple
part of Theorem~\ref{thm:log}; Section~\ref{sec:transfer} proves the
subset part. Notice also the general consequence of \eqref{eq:L2}:
\begin{equation}\label{eq:benchmark}
 k\ge2q+L+3\quad\Longrightarrow\quad\PI(\Fail)\le\frac1{2q}.
\end{equation}

\section{A sharper lattice count and Proposition M}\label{sec:M}
\begin{lemma2prime}
For integers $s,B\ge1$, the number of lattices in $\Z^s$ generated
by vectors of $[-B,B]^s\cap\Z^s$ is at most
\begin{equation}\label{eq:sharper}
 (s+1)2^s(2B+1)^{s^2}(s!B^s)^s.
\end{equation}
\end{lemma2prime}
\begin{proof}
Let $d>0$ be the rank. Choose independent box vectors
$v_1,\ldots,v_d$ in $\Lambda$ and a set $S$ of $d$ coordinates giving
a nonsingular minor $M$. Put $V=\Q\Lambda$ and
$\Lambda_0=\sum_i\Z v_i$. Coordinate projection $\pi$ restricts to
an isomorphism $V\to\Q^d$: it takes the chosen basis to independent
vectors. The finite group
\[
 T_0=\Z^d/\pi(\Lambda_0)
\]
has order $D_0=|\det M|\le d!B^d$ and is generated by $d$ elements.
Each subgroup $U$ of $T_0$ is also generated by at most $d$ elements.
Indeed \eqref{eq:subquot} bounds its number of $p$-primary cyclic
factors by $d$ for every prime. Pad the primary decompositions to $d$
factors and combine across primes by the Chinese remainder theorem.
Thus $U$ has a generating ordered $d$-tuple in $T_0^d$, and there
are at most $D_0^d$ possible $U$.

The data $(d,v_1,\ldots,v_d,S,U)$, with
$U=\pi(\Lambda)/\pi(\Lambda_0)$, determine $\Lambda$. They determine
$V$ and the full inverse image $\pi(\Lambda)$ of $U$, and
\[
 \Lambda=\{x\in V:\ \pi(x)\in\pi(\Lambda)\}.
\]
For if $\pi(x)=\pi(y)$ with $y\in\Lambda$, injectivity on $V$ gives
$x=y$; integrality follows as well. Unrealized data can only overcount.
For fixed $d$ there are at most
$(2B+1)^{sd}2^s(d!B^d)^d$ data, which is at most the same expression
with $d=s$. Sum over $d=0,\ldots,s$, including the unique zero
lattice, to prove \eqref{eq:sharper}.
\end{proof}

For $t\ge3$, $\ell=\lceil\log_2t\rceil$, and $B_t$ from
\eqref{eq:box}, this gives at most $2^{A'_t}$ patterns $\LB(u)$, where
\begin{equation}\label{eq:newA}
 A'_t=t^3(10+2\ell).
\end{equation}
In fact $t!\le2^{t\ell}$, $\log_2B_t\le4t+t\ell$, and
$2B_t+1\le4B_t$, so the logarithm of \eqref{eq:sharper} is at most
\[
 t^3(8+2\ell)+t^2(2+\ell)+t+\log_2(t+1).
\]
For integer $t\ge3$, $\ell\le t$ and $\log_2(t+1)\le t$; the
latter follows by induction from $t+1\le2^t$. The remainder is at
most $t^3+2t^2+2t\le2t^3$, because $t^2\ge2t+2$ for $t\ge3$.
This proves \eqref{eq:newA}.

\begin{propM}
Suppose $q\ge128$, $k$ is an integer, $h=k-q>0$, and integers
$t\ge6$, $1\le a\le t$, with $\ell=\lceil\log_2t\rceil$, satisfy
\begin{align}
 A'_t+aq+t+1&\le tq/16, \tag{P1}\label{eq:P1}\\
 q+2t+2+C_t+3+L&\le ah, \tag{P2}\label{eq:P2}\\
 h&\ge3+L. \tag{P3}\label{eq:P3}
\end{align}
Then
\[
 \PI(\Fail)\le2^{-h}+2^{-q/8}+\frac1{8q}
 \le\frac3{8q}<\frac1{2q}.
\]
\end{propM}
\begin{proof}
Use exactly the event and split of Section~\ref{sec:proof1}.
Replace the low-pattern count by \eqref{eq:newA}. On $E$,
\eqref{eq:P1} gives $|\lo(g)|\le2^{A'_t+aq-tq/16}\le2^{-t-1}$.
The unconditional estimate \eqref{eq:high} is unchanged, since it
uses the binary-pattern count. Markov's inequality and the union
bound therefore give
\[
 \PI(E\cap\Fail)\le2^{q+2t+2+C_t-ah}\le2^{-3-L}=\frac1{8q}
\]
by \eqref{eq:P2}. Combine with \eqref{eq:E}. Condition
\eqref{eq:P3} bounds the first exponential by $1/(8q)$. For real
$q\ge128$, $q/8\ge3+\log_2q$: it holds at $128$ since $16>10$,
and the derivative of the difference is
$1/8-1/(q\ln2)>0$. This bounds the second exponential too.
\end{proof}

The conditions involving $h$ are monotone in $h$. This will give
the ``every integer $k$'' assertion of Theorem~\ref{thm:sqrt}.

\section{Parameters and exact arithmetic for Theorems 2 and 3}
\label{sec:parameters}
We use the following elementary facts for real variables. For $x\ge1$,
$\log_2x\le4\sqrt x$, since
$\ln x=2\ln\sqrt x\le2(\sqrt x-1)\le2\sqrt x$ and
$1/\ln2<2$. The function $2^x/x^2$ increases for $x\ge64$,
as its logarithmic derivative is $\ln2-2/x>0$. The functions
$x/\log_2x$ and $x^{3/5}/\log_2x$ increase for $x\ge e$ and
$x\ge e^{5/3}$, respectively: the signs of their derivatives are
those of $\ln x-1$ and $(3/5)\ln x-1$. Also
$q\ge3+\log_2q$ for $q\ge8$, by checking eight and the positive
derivative of $q-\log_2q$. Finally, for integer $t\ge2$,
\[
 2^{\ell(t)-1}<t\le2^{\ell(t)},\qquad
 \ell(t+1)\le\ell(t)+1,\qquad \ell(t)=\lceil\log_2t\rceil.
\]
The step bound follows from $t+1\le2t$.

For transparency all numerical endpoint comparisons used below are
listed in Table~\ref{tab:constants}. Every entry is an integer
inequality. It can be checked by multiplication alone; the supplied
program also checks it. A positive margin in a row means left side
minus right side for the displayed $\ge$ comparison. Define
\[
 Q_1=512\cdot160^5=53\,687\,091\,200\,000,\qquad
 \QB=49\cdot575^2\cdot30=486\,018\,750.
\]
\begin{table}[ht]
\centering\small
\begin{tabular}{@{}cl@{}}
\toprule
Label & Exact comparison\\
\midrule
A1 & $3400^5\,160^5\ge512\cdot960^5\,161^5$\\
A2 & $3400^5\ge32Q_1$\\
A3 & $2^{64}\ge3402^2\,64^2$\\
B0 & $2^{28}\le\QB<2^{29}$, $\ell(575)=10$\\
B1 & $25(575-23)=24\cdot575$\\
B2 & $200^2\,49^2\,575^2\ge1300^2\cdot56\cdot576^2$\\
B3 & $8400^2\cdot28\ge4\QB$\\
B4 & $2^{64}\ge8402^2\,64^2$\\
B5 & $\QB\ge343^2$\\
B6 & $49\cdot4^9\cdot30>2^{28}$\\
C1 & $9^5\QB^3\ge(40000\cdot29)^5$\\
\bottomrule
\end{tabular}
\caption{Exact endpoints. The margins in B2 and B3 are
$353\,944\,360\,000$ and $31\,605\,000$; the margin in C1 is
$4\,678\,751\,774\,701\,839\,942\,138\,671\,875\,000$.}
\label{tab:constants}
\end{table}

\subsection{Theorem 2}
Let $q\ge64$ and $h=k-q\ge3400q^{4/5}$. If $q<Q_1$, A2 gives
$3400^5\ge32Q_1>32q$, hence $h\ge2q$. It follows that
$k-2q\ge q\ge3+L$, so \eqref{eq:benchmark} applies.

If $q\ge Q_1$, take the largest integer $t$ with $512t^5\le q$.
Then $t\ge160$ and $q<512(t+1)^5$. Put $a=\lfloor t/32\rfloor$;
then $t/40\le a\le t/32$, since
$\lfloor t/32\rfloor>t/32-1\ge t/40$ for $t\ge160$.
The correction guarantees
\begin{equation}\label{eq:A-h}
 h\ge960q/t.
\end{equation}
Indeed maximality, $(t+1)/t\le161/160$, and A1 give
\[
 960^5q<512\cdot960^5(t+1)^5
 \le512\cdot960^5(161/160)^5t^5\le3400^5t^5.
\]
Taking positive fifth roots and rearranging gives
$960q/t<3400q^{4/5}\le h$.

Now $10+2\ell\le8t^3$ (use $\ell\le t$ and $t\ge3$), so
\[
 A'_t\le8t^6\le tq/64,\qquad aq\le tq/32,\qquad
 t+1\le2t\le tq/64.
\]
Their sum proves \eqref{eq:P1}. Next
\[
 C_t+2t+2\le16t^3\le512t^5\le q,\qquad3+L\le q.
\]
Thus the left side of \eqref{eq:P2} is at most $3q$, while
\eqref{eq:A-h} gives $ah\ge24q$. Finally $t^2\le t^5\le512t^5\le q$
implies $t\le\sqrt q$, and
$h\ge960\sqrt q\ge4\sqrt q+3\ge L+3$. This proves
\eqref{eq:P3}. All hypotheses of Proposition M hold, proving the
tuple assertion. The subset assertion follows in Section~\ref{sec:transfer}.

\subsection{Theorem 3(b): the main branch}
For positive integers $t$ define
\[
 f(t)=49t^2(10+2\ell(t)).
\]
It is strictly increasing and unbounded. Suppose $q\ge\QB=f(575)$,
and choose maximal $t$ with $f(t)\le q$. Then $t\ge575$,
$\ell\ge10$, and $q<f(t+1)$. Set $a=\lfloor t/24\rfloor$.
Since an integer has remainder at most $23$ modulo $24$,
\begin{equation}\label{eq:B-a}
 t/25\le(t-23)/24\le a\le t/24;
\end{equation}
the first inequality is exactly $t\ge575$.

For \eqref{eq:P1}, use
\[
 tq/48\ge(49/48)t^3(10+2\ell)=A'_t+A'_t/48\ge A'_t+t+1.
\]
The last step holds for $t\ge3$: $A'_t\ge14t^3$,
$t+1\le4t/3$, and $14t^3/(t+1)\ge(21/2)t^2\ge189/2>48$.
Adding $aq\le tq/24$ proves \eqref{eq:P1}.

For \eqref{eq:P2},
\[
 C_t=4t^2(2+\ell)\le2t^2(10+2\ell)\le2q/49.
\]
Directly from $q\ge f(t)\ge t^2$ we have $t\le\sqrt q$; no
stronger hypothesis on $q$ is invoked. By B5, $\sqrt q\ge343$, so
\[
 2t+5+L\le6\sqrt q+5\le7\sqrt q\le q/49.
\]
The left side of \eqref{eq:P2} is at most $52q/49$. In view of
\eqref{eq:B-a}, it suffices that
\begin{equation}\label{eq:B-h}
 h\ge\frac{1300}{49}\frac qt.
\end{equation}
This also gives \eqref{eq:P3}, since
$h\ge26\sqrt q\ge4\sqrt q+3\ge L+3$.

We deduce \eqref{eq:B-h} from $h\ge200\sqrt{qL}$.
Since $t>2^{\ell-1}$ and $\ell\ge10$,
\[
 q\ge f(t)>49\cdot4^{\ell-1}(10+2\ell)
 \ge\frac{735}{2}\,4^\ell>2^{2\ell+8}.
\]
Thus $L>2\ell+8\ge28$. Maximality of $t$ now gives
\begin{align*}
 q&<49(t+1)^2(10+2\ell(t+1))\\
  &\le49(t+1)^2(12+2\ell)
   <56(t+1)^2L
   \le56(576/575)^2t^2L.
\end{align*}
Here $12+2\ell<L+4<(8/7)L$, using $L>28$. By B2,
\[
 (1300/49)^2q<40000t^2L.
\]
Multiply by $q/t^2$ and take positive square roots to obtain
$(1300/49)q/t<200\sqrt{qL}\le h$. Proposition M proves
$\PI(\Fail)\le3/(8q)$ for every integer $k\ge q+200\sqrt{qL}$.

\subsection{Theorem 3(a): the remaining onset range}
If $64\le q<\QB$, monotonicity and B0 give
\[
 q/L<\QB/\log_2\QB\le\QB/28.
\]
By B3, $8400^2L\ge4q$, so $h\ge8400\sqrt{qL}$ implies
$h\ge2q$. Therefore $k-2q\ge q\ge3+L$, and
\eqref{eq:benchmark} applies. For $q\ge\QB$ use the main branch
and $8400\ge200$. This proves the tuple assertion for every required
integer $k$. The endpoint $q=\QB$ belongs to the main branch, so
the ranges leave no gap.

\subsection{Corollary C}
For $k=\lceil q+3q^{4/5}\rceil$ it suffices to prove
$3q^{4/5}\ge200\sqrt{qL}$, equivalently
$9q^{3/5}\ge40000L$. The ratio $q^{3/5}/\log_2q$ is increasing
for $q\ge\QB$. At $q=\QB$, C1 and $\log_2\QB<29$ imply
\[
 9\QB^{3/5}\ge40000\cdot29>40000\log_2\QB.
\]
Thus the sufficient inequality holds throughout this range. The
ceiling preserves the lower bound on $k$, so Theorem~\ref{thm:sqrt}(b)
gives the tuple conclusion. The required upper bound for the subset
conclusion is verified next.

\section{Transfer from independent sampling to a uniform subset}
\label{sec:transfer}
For independent uniform samples the collision probability is at most
\[
 \binom{k}{2}/N=\frac{k(k-1)}{2N}.
\]
If $k(k-1)\le N$, distinctness therefore has probability at least
$1/2$. For nonnegative integer $k$ and $N\ge2$ this condition also
implies $k\le N$: $k\ge N+1$ would give $k(k-1)\ge N(N+1)>N$.
Conditional on distinctness, the ordered sample is uniform on all
ordered distinct tuples. Forgetting order gives exactly model
\textup{(U)}, and the indexed subset counts become the subset counts
of that set. Balance is invariant under ordering, so
\begin{equation}\label{eq:transfer}
 \PU(\Fail)=\PI(\Fail\mid\text{distinct})
 \le\frac{\PI(\Fail)}{\PI(\text{distinct})}
 \le2\PI(\Fail).
\end{equation}
In particular $k\le2^{q/2}$ is a sufficient upper bound.

For Theorem~\ref{thm:log}, $L\ge20$ and
\[
 k\le q+1000q/L+1\le51q+1\le52q.
\]
For real $x\ge16$, $2^x\ge x^4$: equality holds at sixteen and
the logarithmic derivative $\ln2-4/x$ is positive thereafter. Since
$q\ge2^{20}\ge52$, it follows that
$2^q\ge q^4\ge2704q^2\ge k^2$. This proves both $k\le N$
and the distinctness condition, completing Theorem~\ref{thm:log}.

For Theorem~\ref{thm:power}, $k\le3402q$ because $q^{4/5}\le q$
and $1\le q$. By A3 and monotonicity of $2^q/q^2$ for real
$q\ge64$, $(3402q)^2\le2^q$, so \eqref{eq:transfer} applies.
For Theorem~\ref{thm:sqrt}, the stipulated upper bound on $k$
is sufficient. Its interval is nonempty: $L\le q$ and B4 imply
\[
 \lceil q+8400\sqrt{qL}\rceil\le8402q\le2^{q/2}
 \qquad(q\ge64).
\]
The assertion for constant $200$ follows a fortiori. Finally, in
Corollary C, $k\le5q\le2^{q/2}$, using the same bound with $8402$.
This proves all the subset conclusions. No conclusion for arbitrary
larger $k$ in model \textup{(U)} is needed or asserted.

\section{Scope and limits}\label{sec:limits}
Remark~\ref{rem:scope} gives the finite-size comparison with $L^2$.
The group-uniform statements include elementary abelian groups, cyclic
groups of every order, and mixed-primary groups: $r_G$ takes all primes
dividing $N$ into account. In particular, neither a rational-rank
approximation nor a prime-exponent hypothesis is implicit in the proof.

\begin{proposition}[A limit of the displayed sufficient conditions]
If $q\ge128$, $t\ge6$, $a\ge1$, and $h>0$ satisfy
\eqref{eq:P1}--\eqref{eq:P2}, then
\begin{equation}\label{eq:barrier}
 h>16q/t+64t(2+\ell)
   \ge64\sqrt{q(2+\ell)}\ge128\sqrt q.
\end{equation}
\end{proposition}
\begin{proof}
The positive terms in \eqref{eq:P1} give $a<t/16$.
Equation~\eqref{eq:P2} gives $ah>q+C_t$. Since $a>0$,
$h>(q+C_t)/a>16(q+C_t)/t=16q/t+64t(2+\ell)$.
Apply $x+y\ge2\sqrt{xy}$, and use $\ell\ge2$.
\end{proof}
This is only a limitation of these sufficient conditions. It is not
a lower bound for sharper estimates using realized patterns, their
ranks, their weights, or extra structure of the group.

There is also a much smaller necessary correction for a theorem
uniform over groups in model \textup{(I)}. For $G=(\Z/2)^n$,
$N\ge4$, and $h=k-q\ge1$, let $E_\chi$ be the event that a nonzero
character annihilates every sample. It forces failure, since all sums
then lie in a proper kernel. Each $E_\chi$ has probability $2^{-k}$,
and two distinct nonzero characters have joint kernel of index four,
so $\PI(E_\chi\cap E_\psi)=2^{-2k}$. Bonferroni gives
\begin{align*}
 \PI(\Fail)&\ge(N-1)2^{-k}-\binom{N-1}{2}2^{-2k}\\
 &\ge(N-1)2^{-k}(1-2^{-h}/2)\ge(3/8)2^{-h}.
\end{align*}
The last bound uses $(N-1)/N\ge1/2$ and $1-2^{-h}/2\ge3/4$.
Consequently $\PI(\Fail)\le1/q$ forces
\[
 h\ge\log_2q-\log_2(8/3).
\]
This argument concerns model \textup{(I)} on elementary abelian
two-groups. It is neither a fluctuation claim for every group nor a
lower bound proved here for model \textup{(U)}.

\subsection{Heuristic discussion (not used in the theorems)}
One motivation for a possible square-root scale is the balance between
the spatial factor $2^q$ and the number of relation patterns. A
typical-coefficient heuristic would suggest $|a_\chi|$ cannot all be
much below $2^{-k/2}$ and hence that a low-rank cutoff must have $a$
comparable to at most $t/2$. This heuristic is false for balanced
samples in elementary abelian two-groups, where every nonprincipal
coefficient vanishes. Even in other groups it is not proved here.
Combining the heuristic cutoff with the spatial cost suggests
$h$ of order $q/t$. There are at least $2^{t^2}$ binary-generated
lattices in $\Z^{2t}$: the $2^{t^2}$ graph subspaces
$\{(x,Mx):x\in\F_2^t\}$, as $M$ varies, have distinct mod-two
images, and their binary lifts generate distinct lattices. Paying
for all such patterns suggests a second cost of order $t$ in $h$.
Balancing these two suggested costs gives a square-root scale, while
the bounds actually used here include a logarithmic loss. The count
of all lattices does not establish a lower bound on the number or
weight of the high patterns realized for a fixed group. Thus this
paragraph proves no barrier beyond \eqref{eq:barrier}.

The large constant $8400$ also reflects the chosen hand-over to the
$L^2$ branch. This is a comment about that strategy, not a lower bound
on constants in other parameter choices. For example $f(48)=2\,483\,712$
and the same cutoff would have $a=\lfloor48/24\rfloor=2$, whereas
just below $f(48)$ maximality gives $t\le47$ and $a\le1$. Such
small cutoffs make this particular hand-over costly. No optimized
universal threshold is claimed.

\subsection{Obstructions encountered in simpler approaches}
These observations explain the ingredients used above; none is an
impossibility theorem for alternative proofs.
\begin{enumerate}
\item Direct central moments with a union bound retain the terms
$(N-1)2^{-k}$ from a repeated annihilating character in elementary
two-groups. The additional spatial factor $N$ prevents that estimate
from closing for $h\ll q$. Likewise \eqref{eq:L2} alone gives
$4\,2^{q-h}$, which is uninformative in that regime.
\item Fresh averaging by a random translate gives exact $L^2$
contraction. For nonnegative mean-one $f$, put
$S_p(f)=N^{-1}\sum_g f(g)^p$ and let $T_xf$ denote translation.
Binomial expansion gives
\[
 \E_xS_p((f+T_xf)/2)=2^{1-p}S_p(f)
 +2^{-p}\sum_{j=1}^{p-1}\binom pj S_j(f)S_{p-j}(f).
\]
For $p\ge3$, H\"older interpolation gives
$S_j(f)\le S_{p-1}(f)^{(j-1)/(p-2)}$. Hence the sum term is at
most $(1-2^{1-p})S_{p-1}(f)$. This recurrence still propagates the
second-moment rate $1/2$ and supplied no needed maximum estimate.
\item Taking absolute values before orthogonality loses the exact
binary-relation identity \eqref{eq:moment}. The elementary
Cauchy--Schwarz substitute only gives a square-root relation
probability, losing the required cancellation between $N^r$ and
$2^{-rk}$.
\item A signed remainder cannot be bounded by first-moment Markov.
In $C_5$, four characters with exponents $(1,1,1,2)$ have principal
product, but at a generator the product of their $1+\chi$ factors is
$-16\cos^3(\pi/5)\cos(2\pi/5)<0$. The squared modulus in
\eqref{eq:high} avoids that problem.
\item Binary or signed-binary relations within $u$ need not preserve
rank after combining tuples. If $u_0$ has order five, $(u_0,2u_0)$
has no nonzero relation in $\{-1,0,1\}^2$. But
$(u_0,2u_0,-u_0,-u_0)$ has binary relations
$(1,0,1,0),(1,0,0,1),(0,1,1,1)$; their quotient is infinite cyclic,
and their coordinate intersection contains $(2,-1,0,0)$. The
respective rank parameters are two and one. The enlarged box in
Lemma~\ref{lem:intersection} is essential to this argument.
\item Counting arbitrary subsets of the binary cube costs
$2^{2^{2t}}$. Lemma~\ref{lem:coarse} replaces this by an integer
lattice count. Rational rank alone is also inadequate: $\Z/2\Z$
has rational rank zero but has $N$ homomorphisms into $(\Z/2)^n$.
\item The common coefficient bound alone gives only
$|D(g)|\le(N-1)N^{-1/16}$. In an elementary three-group the expected
Fourier $\ell^1$ norm is $(N-1)(2/3)^k$, also too large near $k=q$.
A split into two sample blocks followed by Cauchy--Schwarz requires
small Fourier $\ell^2$ error in both blocks and did not retain the
leading coefficient one.
\item In primal subset-incidence moments, torsion induces
dependencies beyond repeated subsets, already in elementary
two-groups. An independent-row approximation therefore fails;
\eqref{eq:moment} keeps those dependencies exactly. Higher even
moments of $\hi$ were not developed into a stronger bound: combining
$2m$ tuples gives length $2mt$, but the available rank guarantee is
still only $a$, while the binary-pattern exponent grows with that
length. No correction below the proved square-root logarithmic
scale is established here.
\end{enumerate}

\section{Reproduction and finite checks}\label{sec:reproduce}
The mathematical proof is analytic and uses only the finite integer
comparisons explicitly displayed above. The companion programs are
checks, not a sampling experiment at the enormous onsets and not a
formal proof. They use integer group operations and exact modular
elimination, JavaScript \texttt{BigInt}, or Python integers and
\texttt{fractions.Fraction}; no solver or floating-point acceptance
criterion is used. The two programs are
identical to the reviewed versions except for the removal of an internal output path and of project labels in comments; the archived JSON outputs are unchanged. Their short filenames retain their original check-series
labels; these labels play no mathematical role.

The first program, \texttt{cr19-exact-checks.mjs}, reports $7072$
assertions. For each of
\[
 C_2,C_3,C_4,C_5,C_6,C_7,C_8,C_9,C_{10},C_{12},
 C_2^2,C_2^3,C_3^2,C_2\times C_4,C_2\times C_6,
\]
it enumerates all ordered character tuples of lengths two and four,
including principal characters for pattern checks. It verifies the
cube bound and the count of tuples per pattern. The binary relation
set is an exact key for the generated relation lattice, since the
lattice's intersection with the binary cube is exactly that set.
It checks the fourth-moment relation counts $16,8,6$ at character
orders one, two, and at least three. It also enumerates every ordered
sample of sizes $k=0,1,2,3$ and checks, for $s=2,4$, the identities
\begin{align*}
 \sum_A(NR_A(0)-2^k)^s
   &=N^k\sum_{u\in(H\setminus\{0\})^s}r(u)^k,\\
 \sum_A\sum_g(NR_A(g)-2^k)^s
   &=N^{k+1}\sum_{\substack{u\in(H\setminus\{0\})^s\\\sum_ju_j=0}}r(u)^k,
\end{align*}
where $r(u)=\#\{e\in\{0,1\}^s:\sum_je_ju_j=0\}$.
It checks the three integer bases for the spectral estimate and
induction, \eqref{eq:induction} for $64\le t\le2048$, and
$\lfloor t/32\rfloor\ge t/40$ for $500\le t\le2048$.
These sweeps supplement the general arguments, and do not check
every onset inequality of Theorem~\ref{thm:log} by themselves.

The second program, \texttt{cr19b\_checks.py}, reports $53\,325$
assertions. It checks Table~\ref{tab:constants}, the two onset
parameters, the elementary integer inequalities for $3\le t\le5000$
(factorials through $1000$), and both transfer comparisons for integer
$64\le q\le4096$. It enumerates all box-generated lattices by
closure under adjoining a box vector, with canonical row Hermite
normal forms. The exact totals are
\[
\begin{array}{c|rrrrrr}
 (s,B)&(1,1)&(1,2)&(1,3)&(2,1)&(2,2)&(3,1)\\\hline
 \text{lattices}&2&3&4&7&31&64
\end{array}
\]
and all satisfy Lemma~2$'$. It also evaluates
\eqref{eq:P1}--\eqref{eq:P3} at $q=Q_1m$ for
$m\in\{1,2,7,10^3,10^6\}$, and at $q=\QB m$ for
$m\in\{1,2,3,10,10^3,10^6,10^{12}\}$, using integer logarithm
bounds and integer square roots in the conservative direction. These
sample parameters are checks only; the proofs above cover all real
$q$ in the stated ranges. The minimal positive integer constants
passing the isolated comparisons A1, B2, and B3 are $3364$, $199$,
and $8333$. This is not a claim of optimal theorem constants.

From the bundle directory run
\begin{verbatim}
UV_THREADPOOL_SIZE=1 node --single-threaded code/run-checks.mjs results
\end{verbatim}
with Node.js and Python~3 available. The runner executes scratch copies,
redirects the first program's legacy output destination, and compares
the results with the archived JSON after removing timing fields. It
never changes the archived JSON outputs. The checks are sequential and use
one computation thread. The stored original runtimes are about $0.12$
and $50.34$ CPU seconds; times vary by machine. The second program's
last sample uses a long linear parameter search, which affects runtime
but not the argument. The README gives individual commands and build
instructions. \texttt{SHA256SUMS} records the publication files.

\section*{AI use and verification status}
The mathematical development and manuscript preparation used GPT-6
Astra (OpenAI) and Claude Opus~5.5 (Anthropic). The first bound was
developed using GPT-6 Astra and received two independent Claude proof
reviews with verdict PASS, together with a separate exact-arithmetic
verification. The sharper lattice count and subsequent bounds were
developed using Claude and received two independent GPT-6 Astra
reviews with verdict PASS-WITH-REPAIRS, and a separate numerical
verification with verdict PASS. The repairs replaced an inapplicable
parameter citation by a direct inequality, restricted the methodological
barrier to the sufficient conditions actually proved, and clarified
ceilings and the scope of the lower-bound example. A subsequent delta
verification returned PASS; theorem statements, constants, and onsets
were unchanged. All remaining editorial points are incorporated here.
The exact checks shipped with the paper are the archived author
programs and outputs; reviewer programs are not part of this bundle.
These reviews were performed using AI systems, not human peer review.
The present manuscript was assembled from the reviewed reports and
checked again against their arguments. The literature comparison is
bounded to the search of 12 October 2026 described in
Section~\ref{sec:related}. No Lean or other proof-assistant
formalization of these explicit bounds is supplied.

\begin{thebibliography}{9}
\bibitem{EP1179}
T. F. Bloom, \emph{Erd\H{o}s Problems}, Problem 1179,
\url{https://www.erdosproblems.com/1179}, literature search dated
12 October 2026.
\bibitem{ER65}
P. Erd\H{o}s and A. R\'enyi, \emph{Probabilistic methods in group theory},
J. Analyse Math. \textbf{14} (1965), 127--138.
\url{https://users.renyi.hu/~p_erdos/1965-15.pdf}.
\bibitem{EH76}
P. Erd\H{o}s and R. R. Hall, \emph{Probabilistic methods in group theory II},
Houston J. Math. \textbf{2} (1976), 173--180.
\url{https://users.renyi.hu/~p_erdos/1976-34.pdf}.
\bibitem{Hall77}
R. R. Hall, \emph{Extensions of a theorem of Erd\H{o}s--R\'enyi in
probabilistic group theory}, Houston J. Math. \textbf{3} (1977),
225--234. Bibliographic record only; theorem statements not accessed.
\url{https://www.math.uh.edu/~hjm/vol03-2.html}.
\bibitem{EH78}
P. Erd\H{o}s and R. R. Hall, \emph{Some new results in probabilistic
group theory}, Comment. Math. Helv. \textbf{53} (1978), 448--457.
\url{https://users.renyi.hu/~p_erdos/1978-45.pdf}.
\bibitem{DD24}
M. Dhar and Z. Dvir, \emph{Linear Hashing with $\ell_\infty$ guarantees
and two-sided Kakeya bounds}, TheoretiCS \textbf{3} (2024), article 8;
FOCS 2022 precursor.
\url{https://theoretics.episciences.org/13338/pdf}.
\bibitem{Lund26}
B. Lund, \emph{Two-sided linear hashing and quadratic density bounds
for smooth lattice coverings}, arXiv:2609.20351v1, 17 September 2026.
\url{https://arxiv.org/abs/2609.20351}.
\end{thebibliography}
\end{document}
